% 177-20(177쪽) — 연이율 2 %·1년마다 복리로 매년 초에 a원씩 5년 동안 적립할 때 5년째 말의 적립금 원리합계를 나타낸 그림.
% 스캔 화질이 낮아 TikZ 로 다시 그림. 위쪽 파란 시간축(1년 초~5년 말)·해마다 적립한 a 가 5년째 말까지 불어나는 빨간 화살표·
% 오른쪽 원리합계 칸과 묶음 설명까지 원문 그대로 옮겼다. 원문의 빈칸 상자는 비운 채 두어 답이 그림에서 읽히지 않게 했다.
\begin{tikzpicture}[line width=0.5pt, >=stealth, font=\footnotesize]
  % 위쪽 시간축(1년 초 ~ 5년 말)
  \draw[blue!55, line width=0.6pt] (0,0) -- (5.25,0);
  \draw[blue!55] (0,-0.1) -- (0,0.1);
  \draw[blue!55] (1.05,-0.1) -- (1.05,0.1);
  \draw[blue!55] (2.1,-0.1) -- (2.1,0.1);
  \draw[blue!55] (3.15,-0.1) -- (3.15,0.1);
  \draw[blue!55] (4.2,-0.1) -- (4.2,0.1);
  \draw[blue!55] (5.25,-0.1) -- (5.25,0.1);
  \node[above=0pt, align=center, font=\scriptsize] at (0,0.1) {1\\년\\초};
  \node[above=0pt, align=center, font=\scriptsize] at (1.05,0.1) {2\\년\\초};
  \node[above=0pt, align=center, font=\scriptsize] at (2.1,0.1) {3\\년\\초};
  \node[above=0pt, align=center, font=\scriptsize] at (3.15,0.1) {4\\년\\초};
  \node[above=0pt, align=center, font=\scriptsize] at (4.2,0.1) {5\\년\\초};
  \node[above=0pt, align=center, font=\scriptsize] at (5.25,0.1) {5\\년\\말};
  \node[right=1pt, font=\scriptsize] at (5.25,0) {(단위: 원)};
  % 1년 초에 적립한 a — 5년 동안
  \draw[red!65, ->] (0,-1.15) -- (5.22,-1.15);
  \draw[red!65] (0,-1.26) -- (0,-1.04);
  \draw[red!65] (1.05,-1.26) -- (1.05,-1.04);
  \draw[red!65] (2.1,-1.26) -- (2.1,-1.04);
  \draw[red!65] (3.15,-1.26) -- (3.15,-1.04);
  \draw[red!65] (4.2,-1.26) -- (4.2,-1.04);
  \node[left=1pt] at (0,-1.15) {$a$};
  \node[fill=yellow!45, ellipse, inner sep=1.2pt, font=\scriptsize] at (2.66,-0.8) {5년};
  % 2년 초에 적립한 a — 4년 동안
  \draw[red!65, ->] (1.05,-1.95) -- (5.22,-1.95);
  \draw[red!65] (1.05,-2.06) -- (1.05,-1.84);
  \draw[red!65] (2.1,-2.06) -- (2.1,-1.84);
  \draw[red!65] (3.15,-2.06) -- (3.15,-1.84);
  \draw[red!65] (4.2,-2.06) -- (4.2,-1.84);
  \node[left=1pt] at (1.05,-1.95) {$a$};
  \node[fill=yellow!45, ellipse, inner sep=1.2pt, font=\scriptsize] at (2.66,-1.6) {4년};
  % 가운데 줄임표
  \node at (2.66,-2.72) {$\vdots$};
  % 5년 초에 적립한 a — 1년 동안
  \draw[red!65, ->] (4.2,-3.45) -- (5.22,-3.45);
  \draw[red!65] (4.2,-3.56) -- (4.2,-3.34);
  \node[left=1pt] at (4.2,-3.45) {$a$};
  \node[fill=yellow!45, ellipse, inner sep=1.2pt, font=\scriptsize] at (4.72,-3.1) {1년};
  % 5년째 말의 원리합계 칸
  \draw[line width=0.4pt] (5.3,-0.92) -- (5.3,-3.78);
  \node[anchor=west] at (5.45,-1.15) {$a\times1.02^{5}$};
  \node[draw, line width=0.35pt, minimum width=1.5cm, minimum height=0.48cm, inner sep=0pt, anchor=west] at (5.45,-1.95) {};
  \node at (6.2,-2.72) {$\vdots$};
  \node[draw, line width=0.35pt, minimum width=1.5cm, minimum height=0.48cm, inner sep=0pt, anchor=west] at (5.45,-3.45) {};
  % 오른쪽 묶음 표시와 설명
  \draw[line width=0.35pt] (7.08,-0.92) -- (7.22,-0.92) -- (7.22,-3.78) -- (7.08,-3.78);
  \draw[line width=0.35pt] (7.22,-2.35) -- (7.34,-2.35);
  \node[anchor=west, text width=3.2cm, align=left, font=\footnotesize] at (7.46,-2.35)
    {첫째항이 \blank[2.6em], 공비가 \blank[1.6em]인 등비수열의 첫째항부터 제\blank[1.1em]항까지의 합};
\end{tikzpicture}
